\section{Parity of changes of initial sector identifiers}
\label{sec:area_law_initial_sector_cycle}

Fix an arbitrary origin $o\in\mathbb R^2$, a finite endpoint set
$Z\subset\mathbb Z^2$, a natural width $C\ge2$, and an initial layer
$k_0\ge50{,}000{,}000$. Choose arbitrary valid pitch residues at each
layer. Let $\mathcal I$ be the resulting actual initial identifiers,
write $X_i^0$ for their closed birth regions, and write
$\chi(i)\in\{0,1\}$ for their initial colors.
Let $k\ge k_0$, let $z\in F_{k,\ell_k}$ be an actual fine-cell index,
and let $v$ be one of the nine marks of its defining cell.
Put $\ell=\ell_k-5$ and $r=2^\ell/2=t_k/64$.

Use the all-midpoint fan of the square
$\overline B_\infty(v,r)=v+[-r,r]^2$ and its eight-element index set
$\mathcal J$. Let $\tau$ be the counterclockwise successor permutation
of Definition~\ref{def:area_law_cell_fan_successor}.
Let $\sigma:\mathcal J\to\mathcal I$ be the unique actual assignment
of Theorem~\ref{thm:area_law_actual_sector_assignment}, obtained from
the working fan of half-side $2r=t_k/32$.

% Source: 10-geometry.tex, prop:two-families, lines 299--323,
% and geometry:initial-stars, lines 361--370;
% revision adc7f1241b42e322a6451854ab7e4b4c146bf78a.
% This is the finite parity consequence of the actual adjacent-color
% assertion. Distinct-ray identification and the full star assertion
% are separate. The assignment is derived, not a supplied hypothesis.

\begin{theorem}[Even number of changes of actual initial identifiers]
    \label{thm:area_law_initial_sector_changes_even}
    \lean{TNLean.PEPS.AreaLaw.Geometry.initialRegion_sector_assignment_changes_even}
    \leanok
    \uses{thm:area_law_actual_sector_assignment,
        def:area_law_cell_fan_successor,def:area_law_initial_region_colors}
    The number of changes of actual initial identifier around these
    eight fan members is even:
    \begin{align}
        \left|\{s\in\mathcal J:
        \sigma(s)\ne\sigma(\tau(s))\}\right| & \in2\mathbb N.
    \end{align}
\end{theorem}
\begin{proof}\leanok
    \uses{thm:area_law_initial_sector_colors,
        thm:area_law_cell_fan_successor_endpoints}
    Put $c_s=\chi(\sigma(s))$. The successor endpoint identity makes
    $s$ and $\tau(s)$ adjacent, so the actual adjacent-color assertion
    gives
    \begin{align}
        \sigma(s)\ne\sigma(\tau(s))
         & \quad\Longleftrightarrow\quad c_s\ne c_{\tau(s)}.
    \end{align}
    Let $A=\{s\in\mathcal J:c_s=0\}$ and $B=\tau^{-1}(A)$.
    The binary values of $c_s$ show that the change set is
    $(A\setminus B)\cup(B\setminus A)$, a disjoint union.
    Since $\tau$ is a permutation, $|A|=|B|$, and therefore
    $|A\setminus B|=|B\setminus A|$. With
    $d=|A\setminus B|$, the cardinality of the change set is $d+d$,
    which is even.
\end{proof}
